%! Counting and Probability with Counting — Unit 3, Lesson 3.5 — Lesson Plan
% GRADUAL-RELEASE plan, Unit 2-on timing 5 / 45 / 5 / 5: warm-up / guided notes and practice /
% spoken debrief / close and ASSIGN the homework (not started in class). The guided notes carry
% the release ROW BY ROW in one Main Ideas / Notes table: I do / we do in rows 1-4 (problems
% 1-12, ~25 min) -> we do the Guided Practice row (problems 13-16, ~20 min). Teacher-facing.
\documentclass[10pt]{article}
\usepackage{saar-article}
\usepackage{saar-boxes}
\usepackage{graphicx}   % -article does NOT load graphicx
\graphicspath{{images/}}
\newcommand{\TallMath}[1]{$\displaystyle #1\rule[-1.4em]{0pt}{3.2em}$}
\newcommand{\UnitNumberName}{Unit 3: Categorical Data and Probability \quad}
\newcommand{\LessonNumberName}{Lesson 3.5: Counting and Probability with Counting}
\newcommand{\Pnr}[2]{{}_{#1}P_{#2}}
\newcommand{\Cnr}[2]{{}_{#1}C_{#2}}

\begin{document}
\begin{center}
    {\Large\bfseries \CourseName \\
    {\small\bfseries \UnitNumberName \LessonNumberName}}
\end{center}

\begin{tcolorbox}[colback=frost, colframe=cerulean, boxrule=0.9pt, arc=2mm,
    left=2mm, right=2mm, top=1.5mm, bottom=1.5mm]
    \textbf{Primary Objective:} Students will count outcomes with the Fundamental Counting
    Principle and factorials; calculate permutations and combinations by formula and verify them
    with a calculator; decide which count a situation calls for by checking whether order
    matters, and justify that choice; and find a probability as favorable over total using
    combinations, then predict from it in context.
    \\[2pt]
    \textbf{Standards:} AFDA.DA.3a, g, h, i; A2.ST.3a, b, c, d, e \quad$\cdot$\quad folds in the
    old Lessons 3.7 and 3.8 \quad$\cdot$\quad the last lesson before the Unit~3 review and test
    \\[2pt]
    \textbf{Lesson model:} \emph{Gradual release --- I do, we do, you do.} There is no group
    activity. The notes name the vocabulary as they teach it and deliver the instruction
    \textbf{row by row in one Main Ideas / Notes table} (four rows, problems 1--12: you read
    the definition, mark up the display, and work problems 1, 3, 5, 7, 9; the class works the
    rest); the Guided Practice row (problems 13--16) is worked together with students holding
    the pen. \textbf{The homework is the individual practice} --- assigned at Close \& Assign
    and done outside class. The debrief is a \emph{spoken} cold check and one board item; there
    is no exit ticket.
    \\[2pt]
    \textbf{Conventions on the page:} permutations and combinations count \emph{without
    repetition}, as the standard states; only the Fundamental Counting Principle problems (2 and
    homework item~1a) allow a choice to repeat. Notation is $\Pnr{n}{r}$ and $\Cnr{n}{r}$; the
    calculator screens print \texttt{nPr} and \texttt{nCr}.
\end{tcolorbox}

\renewcommand{\arraystretch}{1.4}
\boxguard
\begin{skillbox}[Learning Targets \& Key Understandings]{goldbox}
\begin{tabularx}{\linewidth}{@{}X|X@{}}
\textbf{Learning targets (student language)} & \textbf{Key understandings (the ``why'')} \\[2pt]
\hline
\vspace{-6pt}
\begin{itemize}[leftmargin=1.1em, nosep, after=\vspace{2pt}]
  \item I can use the \textbf{Fundamental Counting Principle} and \textbf{factorial} notation
        to count outcomes. \hfill \textit{supports AFDA.DA.3h--i}
  \item I can calculate a \textbf{permutation} when order matters and verify it with a
        calculator. \hfill \textit{AFDA.DA.3h; A2.ST.3b, e}
  \item I can calculate a \textbf{combination} when order does not matter, verify it with a
        calculator, and explain why dividing by $r!$ works.
        \hfill \textit{AFDA.DA.3i; A2.ST.3c, e}
  \item I can compare permutations and combinations and choose by checking whether order
        matters. \hfill \textit{AFDA.DA.3g; A2.ST.3a}
  \item I can find a probability as favorable over total using combinations and predict from
        it. \hfill \textit{AFDA.DA.3a; A2.ST.3d}
\end{itemize}
&
\vspace{-6pt}
\begin{itemize}[leftmargin=1.1em, nosep, after=\vspace{2pt}]
  \item \textbf{Target misconception:} using the wrong count because order was not checked ---
        ``both pick $3$ of $10$, so they are the same.'' \textbf{The rule that replaces it:}
        swap two people. If the outcome changes, order matters (permutation); if not, it does
        not (combination). The same $n$ and $r$ give $720$ or $120$: \emph{the question
        decides.}
  \item \textbf{Second misconception:} the word \emph{choose} means combination. Problem 12
        says \emph{choose} and has roles; the swap is the test, not the verb.
  \item Each group of $r$ appears $r!$ times among the permutations, which is why
        $\Cnr{n}{r} = \Pnr{n}{r} / r!$.
  \item A probability with counting uses the same kind of count on top and bottom. Prediction
        is an expected count (AFDA.DA.3a), not a guarantee.
\end{itemize}
\end{tabularx}
\end{skillbox}

\boxguard[24]
\begin{skillbox}[{Vocabulary, Concepts, \& Theorems --- taught directly in the guided notes}]{greenbox}
\small
\renewcommand{\arraystretch}{1.35}
\begin{tabularx}{\linewidth}{@{} >{\bfseries}p{3.4cm} X @{}}
Fundamental Counting Principle & If one decision can be made $n$ ways and another $m$ ways, the
       two together can be made $nm$ ways. \\
Factorial & $n! = n(n-1)(n-2)\cdots(2)(1)$, with $0! = 1$. \\
Permutation & An arrangement of $r$ objects from $n$, no repetition, when order matters:
       \TallMath{\Pnr{n}{r} = \dfrac{n!}{(n-r)!}} \\
Combination & A selection of $r$ objects from $n$, no repetition, when order does not matter:
       \TallMath{\Cnr{n}{r} = \dfrac{n!}{r!\,(n-r)!}} \\
Theoretical probability & $P(E) = \dfrac{\text{favorable outcomes}}{\text{total possible
       outcomes}}$ when outcomes are equally likely (Lesson~3.2). \\
\end{tabularx}
\end{skillbox}

\begin{fixedskillbox}[Lesson at a Glance --- \MeetingLength]{frost}
\renewcommand{\arraystretch}{1.25}
\begin{tabularx}{\linewidth}{@{}l c X X@{}}
\textbf{Phase} & \textbf{Min} & \textbf{Students} & \textbf{Teacher} \\
\hline
Warm-Up & 5 & Evaluate products, cancel, multiply and simplify fractions --- silent &
  Debrief item~2 aloud; leave $\frac{10 \cdot 9 \cdot 8}{3 \cdot 2 \cdot 1} = 120$ on the board
  uninterpreted \\
Guided Notes \& Practice & 45 & Fill the vocabulary box as each term is named; work rows 1--4
  (problems 1--12) with the pen in their hand; then the Guided Practice row, \emph{Priya's
  Playlist} (13--16), together & \textbf{I do / we do, row by row}: read the definition, mark
  up the display, work problems 1, 3, 5, 7, 9; the class works the rest (25) $\to$ \textbf{we
  do} Guided Practice 13--16 (20) \\
Debrief & 5 & Pens down --- answer aloud, nothing to fill in & Put problem~15 back on the
  board; ask the cold check \\
Close \& Assign & 5 & Write down the homework, the due date, and item~4 as the first item &
  Announce the homework and its form (DeltaMath + packet items 4, 5, 6), preview the Unit~3
  review and test \\
\end{tabularx}
\end{fixedskillbox}

\begin{fixedskillbox}[Warm-Up --- Activate Prior Knowledge (5 min)]{frost}
\begin{tabularx}{\linewidth}{@{}X X@{}}
\begin{minipage}[t]{\linewidth}\vspace{0pt}
    \textbf{The three items and what each seeds}
    \begin{itemize}[leftmargin=1.1em, itemsep=1pt, topsep=2pt]
      \item \textbf{1.} $5 \cdot 4 \cdot 3 \cdot 2 \cdot 1 = 120$; $10 \cdot 9 \cdot 8 = 720$;
            $3 \cdot 2 \cdot 1 = 6$. \textbf{Seeds:} problems 3, 4, and 9, and the $6$ of
            problem 11.
      \item \textbf{2.} Cancelling gives $35$ and $120$. \textbf{Seeds:} problems 7 and 10.
      \item \textbf{3.} $\frac{5}{12} \cdot \frac{4}{11} \cdot \frac{3}{10} = \frac{60}{1320} =
            \frac{1}{22}$ and $\frac{10}{220} = \frac{1}{22}$. \textbf{Seeds:} problem 15, and
            reaches back to the Lesson~3.3 multiplication rule.
    \end{itemize}
\end{minipage}
&
\begin{minipage}[t]{\linewidth}\vspace{0pt}
    \textbf{Running it}
    \begin{itemize}[leftmargin=1.1em, itemsep=1pt, topsep=2pt]
      \item \textbf{Debrief item~2 aloud} on its right-hand fraction, and leave it on the
            board. Do not say what $120$ counts; problem~10 will.
      \item Item~3's noticing is the tell for arithmetic: students who multiplied everything
            out and could not simplify need cancelling modeled at problem~7.
      \item No calculators on the warm-up; the calculator arrives in row~2 as a pre-made
            screen.
    \end{itemize}
\end{minipage}
\end{tabularx}
\end{fixedskillbox}

\boxguard
\begin{skillbox}[Hook --- before anyone sits down]{frost}
The hook is on \textbf{page 1 of the notes}, under the vocabulary box, and on the board:
\textit{``Picking $3$ of $10$ students for a committee and picking a president, vice president,
and treasurer from $10$ students give the same number of possibilities.''} Students circle
\emph{True} or \emph{False} and write a reason \emph{before} anyone speaks; then take the hand
vote and write the split on the board. Many rooms vote \emph{True}: both pick three people from
ten. \textbf{Do not settle it.} Say only that the class votes again at problem~10, in the last
instruction row. Rows 2 and~3 deliberately use other numbers ($8$, $6$, $7$) so that the hook's
$10$ and $3$ do not appear until row~4.
\end{skillbox}

\boxguard
\begin{skillbox}[Guided Notes \& Practice --- I do, then we do (45 min)]{frost}
\textbf{This is the whole lesson.} There is no group activity, and there is no independent
practice set on this page: \textbf{the homework is the individual practice}. The notes are
\textbf{one Main Ideas / Notes table}: four instruction rows (problems 1--12), then the Guided
Practice row (13--16). \textbf{The rhythm in every row is the same}: read the definition aloud,
mark up the display, work the first problem of the grid yourself, then hand the rest to the
room. The vocabulary box on page~1 is filled \emph{as each term is named}.
\begin{multicols}{2}
\textbf{I do / we do, row by row (about 25 min).}

\emph{Row 1, Counting Principle (problems 1--4).} Trace the pre-drawn tree: three fundraisers,
two days each, six paths, so $3 \times 2$. Name the \emph{Fundamental Counting Principle}.
Work problem~1 ($24$ plans) yourself; the class does 2 ($1000$ codes --- repeats allowed,
which is why this principle and not a permutation). Name \emph{factorial}; work problem~3
($5! = 120$), the class does 4 ($10!/7! = 720$, the warm-up's product). Six minutes.

\emph{Row 2, Permutations (problems 5--6).} Read the definition and formula, point to the two
slots ($8$ then $7$), work problem~5 by formula and check it against the top line of the screen
($56$). The class does 6 ($360$ ordered lineups), checking the second line. The screen is
verification, not the method. Six minutes.

\emph{Row 3, Combinations (problems 7--8).} Students fill the display: one committee, $3! = 6$
orderings. Read the formula; work problem~7 ($\Cnr{7}{3} = 35$, the warm-up's cancellation)
and read the screen down: $210$, $210/3!$, $35$. The class does 8, the reason for dividing by
$r!$. Six minutes.

\emph{Row 4, Which Count? (problems 9--12) --- the crux.} The hook's numbers return, side by
side. Students label the two panels. Read the boxed sentence. Work problem~9 live: swap two
officers, the outcome changes, $\Pnr{10}{3} = 720$. \textbf{Re-take the hook vote before
anyone counts problem~10}, then count: $\Cnr{10}{3} = 120$. \textbf{Problem~10 is the crux}.
Problem~11 says why the ratio is exactly $3! = 6$. \textbf{Problem~12 is the trap}: Ty's
\emph{choose} with two different jobs is $90$, not $45$. Seven minutes.

\columnbreak
\textbf{We do (about 20 min).} The Guided Practice row, \emph{Priya's Playlist} (problems
13--16) --- a fresh context, and the probability half of the lesson. \textbf{Students hold the
pen; you hold the questions.} Drive the calculator keystrokes aloud here if students have
them. Problem~13 re-asks the crux on new ground: \emph{does order matter in a mix?} It does
not, so $\Cnr{12}{3} = 220$. Problem~14: $\Cnr{5}{3} = 10$. Problem~15: $\frac{10}{220} =
\frac{1}{22} \approx 0.045$, the warm-up's fraction. Problem~16: $110 \cdot \frac{1}{22} = 5$
days, a prediction and not a promise.

Circulate during Guided Practice and demand the reason, not the number:
\begin{itemize}[leftmargin=1.1em, nosep]
  \item \textbf{Problem 13:} ``Swap two songs in the mix. Is it a different mix?''
  \item \textbf{Problem 14:} ``Why $5$ on top of the combination and not $12$?'' (Only Luna
        Vale's songs are favorable.)
  \item \textbf{Problem 15:} ``Which count is on top? Which is on the bottom? Are they the same
        kind?''
  \item \textbf{Problem 16:} ``Does $5$ mean exactly $5$ days?'' (About $5$; it is a
        prediction.)
\end{itemize}
While circulating, pick the one paper to put up in the debrief --- a problem~15 with the
counts labeled favorable and total. \textbf{This circulation is your formative read.}
\end{multicols}
\end{skillbox}

\boxguard[24]
\begin{skillbox}[Debrief --- whole class, spoken (5 min)]{frost}
Pens down, papers up. \textbf{This debrief is spoken} --- it lives on the board and in the room.
\begin{multicols}{2}
\begin{enumerate}[leftmargin=1.3em, nosep, itemsep=3pt]
  \item \textbf{The one board item: problem~15.} Put up the paper you picked. Ask what goes
        wrong if the top is a combination and the bottom a permutation
        ($\frac{10}{1320}$ is too small by a factor of $6$). Make a student say the sentence:
        \emph{check order before you count.}
\end{enumerate}
\columnbreak
\textbf{Check for understanding --- ask it cold, whole class.} \emph{``Eight runners race. One
question: how many ways to award gold, silver, and bronze? Another: how many ways can $3$ of
the $8$ advance to the finals? Same number? Which is larger, and by what factor?''} Hands down,
thirty seconds of think time, then cold-call three students. Correct answer: \textbf{no}. The
medals have roles, so $\Pnr{8}{3} = 336$; advancing has none, so $\Cnr{8}{3} = 56$. The medals
count is larger by a factor of $3! = 6$. \emph{This replaces the exit ticket.}

\textbf{Formative read.} Sort the answers into three piles: \textbf{(a)} said \emph{no}, named
order or roles, and gave the factor $6$; \textbf{(b)} said \emph{no} but gave no reason beyond
the numbers; \textbf{(c)} said \emph{the same, both are $3$ of $8$} --- order not checked.
\textbf{If pile (c) runs past a third of the room, open the Unit~3 review by re-running row~4
(problems 9--10)}, and say so plainly.

\textbf{Five minutes: the board item and the cold check only.}
\end{multicols}
\end{skillbox}

\boxguard
\begin{skillbox}[Homework --- scored 1 to 5, assigned today, due after two study halls]{goldbox}
Six items on \textbf{locker codes and a fantasy-team draft}: \textbf{1} codes with repeats
($10{,}000$) and without ($5{,}040$) --- the core procedure; \textbf{2} and \textbf{3} are the
\textbf{contrast pair}, the same $15$ players and $3$ picks: ranked spots, permutation $2{,}730$;
bench spots alike, combination $455$; \textbf{4} Devon's ``the answers must be the same'' ---
\emph{the crux head-on}, and the justification (rubric); \textbf{5} the probability of an
all-rookie team, $\frac{20}{455} = \frac{4}{91} \approx 0.044$, and the interpretation (about
$4$ of $91$ teams; rubric); \textbf{6} \emph{the spiral item}, reaching back to Lesson~3.4: two
players drawn without replacement, $\frac{6}{15}$ against $\frac{5}{14}$, so the events are
dependent.

\textbf{Assigned at Close \& Assign and done outside class.} \textbf{Scoring:} 12 points ---
2 per turned-in item, the DeltaMath score scaled to the rest, scored 1 to 5 by the grading
report. \textbf{Item~4 is where the misconception shows.}

\textbf{Packet, Desmos, or DeltaMath --- decided here: DeltaMath + justifications.} The packet
pages are authored and are the default; no Desmos activity asks for the written reasons this
lesson exists for. Items 1--3 are counting procedure --- codes with and without repeats, then
naming and computing a permutation and a combination --- which DeltaMath drills; items 4, 5,
and~6 are the crux justification (why $2{,}730 = 6 \times 455$), a probability with its
interpretation, and the Lesson~3.4 spiral decided by two probabilities, which it cannot grade.
So assign \textbf{DeltaMath} --- Fundamental Counting Principle (L2) $\cdot$ Permutations $\cdot$
Permutations vs. Combinations (One Step) $\cdot$ Simplify Factorial Expressions $\cdot$ Compound
Probability --- for the computation (items 1--3 stay in the packet as an unboxed worked
reference), and students \textbf{turn in packet items 4, 5, and~6}. Still 12 points: 2 per
turned-in item, the DeltaMath score scaled to the other 6. Say so aloud at Close \& Assign.

\textbf{Preview of the next class.} The Unit~3 review: every tool from 3.1--3.5, then the
Unit~3 test. Item~4 is test-style --- say so.

\textbf{Connections \& Big Ideas.} \textit{Unit 3 core idea:} probability describes chance with
numbers, and those numbers depend on correctly counting and classifying what can happen.
\textit{Standards covered:} AFDA.DA.3a (make predictions from theoretical probability),
AFDA.DA.3g (compare permutations and combinations), AFDA.DA.3h and 3i (calculate them without
repetition), A2.ST.3a--c (compare, calculate), A2.ST.3d (use them in contextual problems),
A2.ST.3e (calculate and verify with technology). \textit{Spirals forward to:} the Unit~3
review and test. \textit{Reaches back to:} Lesson~3.2's theoretical probability, Lesson~3.3's
multiplication rule, and Lesson~3.4's conditional probability and dependent events.
\end{skillbox}

\boxguard
\begin{skillbox}[Watch For (while circulating)]{redbox}
\begin{itemize}[leftmargin=1.2em, nosep]
  \item \textbf{Order not checked} (the hook, problems 9--10 and~13, HW~2--4 --- the target
        misconception). Probe: \emph{``Swap two people. Is it a different outcome? Then which
        count is it?''}
  \item \textbf{``Choose means combination''} (problem~12, HW~2). Probe: \emph{``Do the two
        chosen people have the same job?''}
  \item \textbf{Permutation when repeats are allowed} (problem~2, HW~1). Probe: \emph{``After
        one digit is used, how many digits can the next one be?''}
  \item \textbf{Leaving off $r!$ in the bottom} (problems 7 and~10, HW~3). Probe: \emph{``How
        many times did you count this one committee?''}
  \item \textbf{Mixed counts in a fraction} (problem~15, HW~5). Probe: \emph{``Are your top
        and bottom the same kind of count?''}
  \item \textbf{Reading a prediction as a guarantee} (problem~16, HW~5). Probe: \emph{``Could
        it be $4$ days? Could it be $7$?''}
  \item \textbf{Treating a draw without replacement as unchanged} (HW~6). Probe: \emph{``How
        many rookies are left, and out of how many players?''}
\end{itemize}
Cold-call prompts: \emph{``Give me a situation at this school where order matters, and one
where it does not.''} \quad \emph{``Why can a combination never be larger than the permutation
with the same numbers?''}
\end{skillbox}

\boxguard
\begin{skillbox}[Close \& Assign (5 min)]{goldbox}
\textbf{The homework is assigned here and done outside class --- it is not started in class.}
Hand the launch: six items on locker codes and a fantasy-team draft, on the packet pages,
scored 1 to 5, due at the start of the first class after two study halls. Say the form:
\textbf{DeltaMath} for the counting, permutation, and combination practice, and \textbf{packet
items 4, 5, and~6 turned in}. Point out that
the \emph{Keep in Mind} box on the cover and rows 2--4 of the notes are the pages to flip back
to. \textbf{Tell students to do item~4 first if they are short on time}, and that items 2 and~3
are a pair: read both before answering either. Then close on what changed: \emph{``You walked
in able to multiply and cancel. You walk out able to count outcomes, and knowing that the same
two numbers give different answers when the question changes: check order before you count.''}
\textbf{Preview:} the Unit~3 review, then the test.
\end{skillbox}

% ── Teacher notes — one per component, in packet order (three) ──────────────
\begin{teachernote}[Warm-Up]
Five minutes, silent, no calculator. Every number returns: $120$ at problem~3, $720$ at problems 4
and~9, $6$ at problems 8 and~11, $35$ at problem~7, $120$ again at problem~10, and $\frac{1}{22}$
at problem~15. Say so, so students keep the page. \textbf{Do not interpret} the fractions in
item~2: the answer to ``what does $120$ count?'' is the reveal at problem~10. Item~3's two equal
results connect the multiplication rule to today's combinations; say only that the class
will see that fraction again.
\end{teachernote}

\begin{teachernote}[Guided Notes \& Practice]
\textbf{This one document is the lesson --- pace it 25 / 20, then 5 for the debrief and 5 to
assign the homework.} It is one table, and every row runs the same way: you read the
definition, mark up the display, and work the first problem of the grid (1, 3, 5, 7, 9); the
class works the rest with the pen in their hand. Rows 1--3 must move --- six minutes each.

\textbf{Spend the time in row~4, \emph{Which Count?}}, which carries the misconception.
Work problem~9 live, then \textbf{re-take the hook vote before problem~10 is counted.} The trap
(problem~12) comes after the crux on purpose: students who have just learned the swap test are
exactly the ones tempted to go back to reading the verb.

\textbf{Guided Practice (13--16) is where the pen changes hands.} Problem~13 is the order check
in a new context; 14--15 are favorable over total, and 16 is the prediction (AFDA.DA.3a). A
student who uses permutations on both top and bottom ($\frac{60}{1320}$) gets the same
probability; that is correct, and worth saying, but today's habit is combinations when order
does not matter. Mixing them is the error. If time is short, protect 13 and 15.

\textbf{Scope.} The standard says permutations and combinations \emph{without repetition}, so
no permutation or combination on these pages allows a repeat; problem~2 is a Counting Principle
problem, and it says repeats are allowed. Circular arrangements, identical objects, and ``at
least one'' counts are left out.
\end{teachernote}

\begin{teachernote}[Homework]
\textbf{Scored 1 to 5 from 12 points (2 per turned-in item, DeltaMath scaled to the rest), due
at the start of the first class after two study halls; assigned at the close, not started in
class.} The turned-in items keep their boxes. Item~6 is short boxes, scored by value or word
($0.5 + 0.5 + 1$). Items 4 and~5b are rubric boxes (\texttt{rubrics/}): item~4, 1 point for
naming order or roles, 1 for the factor $3! = 6$; item~5b, 1 point for about $4$ of $91$ in
context, unusual but possible. Item~5a takes $0.044$ within $0.001$.

\textbf{Item~4 is the one to read first.} Sort it into three piles --- named order and the
factor (ready); named order only (ready, needs the reason for $6$); agreed with Devon or
compared only the numbers (the misconception survived). If more than a third land in the last
pile, open the Unit~3 review with problems 9--10 before anything else.

\textbf{DeltaMath + justifications is the call for 3.5.} DeltaMath (exact topic names):
Fundamental Counting Principle (L2) $\cdot$ Permutations $\cdot$ Permutations vs. Combinations
(One Step) $\cdot$ Simplify Factorial Expressions $\cdot$ Compound Probability. Students turn in
packet items 4, 5, and~6. Score the three turned-in items as above (item~4 first, then 5, then
6); scale the DeltaMath score to the other 6 points. The items 2--3 contrast-pair check moves to
Permutations vs.\ Combinations (One Step): a student who misses the order-matters problems there
is not checking order --- and item~4 still asks for that reason in writing. Simplify Factorial
Expressions is the $\Cnr{n}{r} = \Pnr{n}{r}/r!$ cancelling, drilled. Collect the turned-in items
at the start of the first class after two study halls and use item~4's piles to open the Unit~3
review.
\end{teachernote}
\end{document}
